\documentclass[11pt]{article}

\usepackage[UTF8]{ctex}
\usepackage{amsmath,amssymb,amsfonts,amsthm}
\usepackage{booktabs,array}
\usepackage{xcolor}
\usepackage{hyperref}

\definecolor{studyblue}{RGB}{35,90,145}
\definecolor{studyteal}{RGB}{25,115,105}
\definecolor{studyorange}{RGB}{185,95,20}

\newtheorem{theorem}{定理}[section]
\newtheorem{lemma}[theorem]{引理}
\newtheorem{proposition}[theorem]{命题}
\theoremstyle{remark}
\newtheorem{remark}[theorem]{注}

\newenvironment{prerequisite}[1]
  {\par\medskip\noindent\color{studyteal}\textbf{基础补充：#1}\par\smallskip\color{black}}
  {\par\medskip}
\newenvironment{calculation}[1]
  {\par\medskip\noindent\color{studyblue}\textbf{计算补充：#1}\par\smallskip\color{black}}
  {\par\medskip}
\newenvironment{checkpoint}[1]
  {\par\medskip\noindent\color{studyorange}\textbf{学习检查：#1}\par\smallskip\color{black}}
  {\par\medskip}

\title{数学论文中文学习版}
\author{}
\date{}

\begin{document}
\maketitle

\begin{center}
\fcolorbox{studyblue}{blue!4}{%
  \parbox{0.9\linewidth}{\small
  \textbf{来源与状态。}
  本学习版依据指定的源论文快照编写。这里填写源文件、版本日期、数学证据等级，
  并说明新增解释和计算是否经过独立核验。}}
\end{center}

\tableofcontents

\section*{快速学习指南}
\addcontentsline{toc}{section}{快速学习指南}

用一段话写出全文的核心机制，再给出分遍阅读计划、证明依赖关系和少量关键量。

\begin{checkpoint}{读正文前}
列出读者开始第一遍前应能回答的两至四个实质问题。
\end{checkpoint}

\section{问题、主结果与证明主线}

准确说明对象、假设、维数、量词、目标量和主定理，然后解释各证明模块之间的关系。

\section{第一个证明模块}

\begin{prerequisite}{本节最低限度的背景}
解释马上要使用的定义、标准公式及其几何意义。
\end{prerequisite}

这里放与源论文对应的定义、命题和证明主干。

\begin{calculation}{决定结论的一步}
逐行展开变量、求导对象、符号、缩放或边界项，并在结尾说明该计算解决了什么障碍。
\end{calculation}

\begin{checkpoint}{本模块}
要求读者不看正文重建本节的输入、关键估计和输出。
\end{checkpoint}

\section*{证明重构清单}
\addcontentsline{toc}{section}{证明重构清单}

按数学依赖列出从精确目标到最终闭合的步骤，并注明最长技术核查的位置。

\begin{checkpoint}{最终掌握标准}
写出读者合上稿件后应能口头解释和在纸上推导的内容。
\end{checkpoint}

\end{document}
